\documentclass[10pt,a4paper]{amsart}
\usepackage[margin=19mm]{geometry}
\usepackage{amsmath,amssymb,mathtools}
\usepackage[scr=rsfs,cal=euler]{mathalfa}
\usepackage{xfrac}
\usepackage[unicode,hidelinks,bookmarks=false]{hyperref}
\newcommand{\ie}{i.e.\ }
\newcommand{\cf}{cf.\ }
\newcommand{\resp}{respectively\ }
\newcommand{\Subsection}{\subsection}
\providecommand{\nobreakdash}{\nobreak\hbox{-}}
\setlength{\emergencystretch}{2em}
\multlinegap0pt
\usepackage{bm}
\usepackage{bbm}
\usepackage{dsfont}
\usepackage{xspace}
\usepackage{cancel}
\usepackage{mathtools}
\usepackage{amsthm}
\makeatletter
\@ifundefined{thm}{\newtheorem{thm}{Thm}}{}
\@ifundefined{lem}{\newtheorem{lem}[thm]{Lemma}}{}
\makeatother
\title[Abelian p-groups with a fixed elementary subgroup or with a ]{Abelian p-groups with a fixed elementary subgroup or with a fixed elementary quotient}
\author{Justyna Kosakowska, Markus Schmidmeier, Martin Schreiner}
\date{Source: arXiv:2312.01451v2 (CC BY 4.0)}
\begin{document}
\maketitle

\noindent\emph{Source and licence.} Abelian p-groups with a fixed elementary subgroup or with a fixed elementary quotient. Version arXiv:2312.01451v2 (CC BY 4.0), distributed under the Creative Commons Attribution 4.0 International licence, as recorded in the arXiv licence metadata for this exact version and shown on the abstract page https://arxiv.org/abs/2312.01451v2. Full rights notice: licenses/arxiv-2312.01451-CC-BY-4.0.txt.\par\smallskip

\noindent\textit{Editorial scope.} Counted unit: the Lem whose source label is \texttt{lem-GT} in the article (source0.tex lines 490--498, source label \texttt{lem-GT}), delivered together with the article's own complete original proof of it (source0.tex lines 500--537, one \texttt{proof} environment of 38 lines). No mathematics is written, repaired or re-derived: statement and proof are the article's own text line for line. Required context retained uncounted: none. Every definition, display and label of those lines is retained unchanged. Omitted from this package and not redistributed: 1--489, 499--499, 538--779. Nothing inside a retained excerpt is omitted or altered: every retained body is delivered exactly as the article's lines are written, so no omission receipt of any kind accompanies this package. Citations: the retained excerpts carry no citation command at all, so no bibliography entry of the article is transcribed and no reference is redistributed. Reference adaptation: none was needed and none was made; every reference inside the delivered unit resolves inside it, so no unresolved reference or citation is produced. Printed slips kept literal: no printed word, symbol, hypothesis or number of the counted statement or of its proof was altered. Layout and class: the article's own class is \texttt{amsart} with packages including amsmath; the wrapper is the lane's pinned \texttt{amsart} editorial wrapper at 10pt on A4, which restates the theorem-like environments the retained text needs and adds the article's own macro definitions that the retained text needs (none), with extra wrapper packages (bm, bbm, dsfont, xspace, cancel, mathtools). The delivered numbering is the unit's own. Assets: no retained span contains a figure environment, a graphics-inclusion command, a PDF-page inclusion, a table or a TikZ picture; no image file is copied into this package, the package declares no asset dependency and no image file is redistributed with it. Licence: version arXiv:2312.01451v2 of this article is distributed under the Creative Commons Attribution 4.0 International licence, as recorded in the arXiv licence metadata for this exact version and shown on the abstract page https://arxiv.org/abs/2312.01451v2. A full rights notice is delivered at licenses/arxiv-2312.01451-CC-BY-4.0.txt. Characters: every retained excerpt is pure ASCII, so the delivered engine, XeLaTeX with its default Latin Modern text font, reports no missing character.
\begin{lem}
  Suppose $T$ is a direct sum of copies of $R/(p^n)$.
  Let $B=\{b_i:i\in I\}$ be a linearly independent set in $T/pT$
  and $C=\{c_i:i\in I\}$ any set of liftings of $B$ to $T$ satisfying $b_i=c_i+pT$.
  Then the submodule $M$ of $T$
  generated by the elements $c_i$ is the direct sum
  $$M=\bigoplus_{i\in I}c_i \mathbb Z.$$
  Moreover, if $B$ is a basis for $T/pT$, then $T=\bigoplus_{i\in I} c_i\mathbb Z$.
\label{lem-GT}\end{lem}

\begin{proof}
  We show by induction that for each $m=n-1,\ldots,0$ the following statement holds:

  $$(P_m) \quad p^mM=\bigoplus_{i\in I} p^mc_iR; \quad\text{if $B$ is a basis then}\quad
  p^mT=p^mM.$$

  We will use that the map
  $$\varphi_m: \frac{T}{pT}\to\frac{p^mT}{p^{m+1}T}, \;x+pT\mapsto p^mx+p^{m+1}T$$
  is an isomorphism of vector spaces.  In particular, the statement $(P_m)$ holds
  for $m=n-1$.


  We deal with the induction step $m+1\to m$.
  Let $B'=\{b_i'=\varphi_m(b_i)  :i\in I\}$.
  Then $B'$ is a linearly independent set and, if $B$ is a basis, then so is $B'$.
  The elements $b_i'$ satisfy $b_i'=p^mc_i+p^{m+1}T$ for $i\in I$.


  The elements $p^mc_i$ form a generating set for $p^mM=\sum_{i\in I}p^mc_iR$.


  To show that the sum is a direct sum, suppose $J\subseteq I$ is a~finite subset
  and $0=\sum_{i\in J} a_ip^mc_i$ where
  $a_i\in R$.   Then in $p^mT/p^{m+1}T$, we have
  $0=\sum a_ib_i'$.  Since  $B'$ is a linearly independent set,
  for all $i\in J$ there exist $a_i'\in R$ such that $a_i=a_i'p$.
  Since $\sum a_i'p^{m+1}c_i=0$ we can use the induction hypothesis $(P_{m+1})$ to conclude
  that $a_i'p^{m+1}c_i=0$ for all $i\in J$, hence each $a_ip^mc_i=0$, so
  the sum is a direct sum.


  It remains to show that if $B$ is a basis then $p^mT=p^mM$.
  Let $x\in p^mT$, then $x+p^{m+1}T=\sum a_ib_i'$ with $a_i\in R$.
  Hence $x-\sum a_ip^mc_i$ is in $p^{m+1}T=p^{m+1}M$, by induction hypothesis $(P_{m+1})$,
  so there are
  $a_i'\in R$ with $x-\sum a_ip^mc_i=\sum a_i'p^{m+1}c_i$.
  This shows that $x\in p^mM$.
\end{proof}

\end{document}
